% !TEX root = ../main.tex
% !TEX program = pdflatex
\newpage
\section{Заключение}

В лабораторной работе реализована программа на языке C для исследования кэш-памяти Apple~M2. Задержка измерялась обходом случайной зависимой цепочки; последовательный обход использовался отдельно для оценки эффективной скорости обработки массива. Каждый замер предварялся двумя прогревочными проходами, затем выполнялись пять измерений. Пользователь провёл три полных запуска; диапазоны ниже показывают разброс медиан между запусками. Полученные выводы сведены в таблицу~\ref{tab:result}.

\begin{table}[H]
\centering
\small
\caption{Итоговые оценки характеристик памяти Apple M2}\label{tab:result}
\begin{tabularx}{\textwidth}{lXX}
\toprule
Уровень & Сведения о кэше & Результат измерения \\
\midrule
L1 данных P-ядер & 128 КиБ по ОС; строка около 64 Б; ассоциативность около 8 каналов & Плато около 1,14 нс; переход около 128--160 КиБ \\
L2 P-кластера & 16 МиБ по ОС; 128 Б сообщает ОС; ассоциативность не определена & Плато около 5,4--5,6 нс; переход растянут примерно на 8--24 МиБ \\
Нижняя часть иерархии & Точная граница этим опытом не устанавливается & 102--106 нс при рабочем наборе 256 МиБ \\
\bottomrule
\end{tabularx}
\end{table}

В размерном sweep задержка около 1,14~нс сохраняется до 128~КиБ, после чего возрастает; это согласуется с размером L1 данных производительного ядра. На больших наборах наблюдается плато около 5,4--5,6~нс и дальнейший рост в диапазоне примерно 8--24~МиБ. macOS сообщает L2 объёмом 16~МиБ на кластер производительных ядер, но эксперимент не определяет его точную границу из-за постепенного перехода, влияния TLB и общего кэша SLC.

В тесте строки зависимая пара чтений показывает повторяемый скачок при переходе от смещения 56 к 64~байтам для набора, размещающегося в L2. Это поддерживает оценку строки L1 около 64~байт. Для набора 256~МиБ рост также заметен около 128~байт и согласуется со значением \code{hw.cachelinesize}, однако этот опыт не изолирует строку L2 от эффектов нижних уровней памяти.

В конфликтном тесте L1 задержка начинает устойчиво расти при добавлении девятой строки; по трём полным запускам медиана при $N=9$ составила 1,76--1,98~нс против примерно 1,14--1,20~нс для $N=4\ldots8$. Это согласуется с 8-канальной ассоциативностью L1. Ассоциативность L2 достоверно не определена: пользовательская программа не контролирует физические адреса и отображение строк в наборы.

При последовательной обработке эффективное отношение размера массива ко времени прохода составило около 40,7--41,1~ГБ/с для шага 64~байта и 33,0--33,1~ГБ/с для шага 128~байт. Это характеристика выбранного цикла, а не предельная пропускная способность DRAM. В целом эксперимент показывает, что наблюдаемая задержка зависимого доступа меняется примерно от 1,14~нс на малом рабочем наборе до 102--106~нс на 256~МиБ, а аппаратные выводы необходимо формулировать с учётом ограничений пользовательского измерения.
